\section{Weighted Chatterjee's \texorpdfstring{$\xi$}{xi}}
\label{sec:weighted-chatterjee}

\subsection{Population coefficient and weighted estimator}

Chatterjee's coefficient is directional: it measures the extent to which the
response $Y$ is determined by the predictor $X$. If $Y$ is not almost surely
constant and $\mu_Y$ denotes its law, the population coefficient is
\begin{align*}
  \xi(X,Y)
  &=
  \frac{
    \displaystyle
    \int
    \operatorname{Var}\!\left\{
      \mathbb{P}(Y\geq t\mid X)
    \right\}\,\mathrm{d}\mu_Y(t)
  }{
    \displaystyle
    \int
    \operatorname{Var}\!\left\{
      \mathbf{1}(Y\geq t)
    \right\}\,\mathrm{d}\mu_Y(t)
  }.
\end{align*}
It belongs to $[0,1]$, equals zero if and only if $X$ and $Y$ are
independent, and equals one if and only if $Y$ is almost surely a measurable
function of $X$ \citep{chatterjee2021}. In general,
$\xi(X,Y)\neq\xi(Y,X)$.

The sample coefficient of \citet{chatterjee2021} orders the observations by
the predictor and measures the total variation of the response ranks along
this path: a response that is a function of the predictor produces a smooth
path, whereas independence produces large jumps. A case-weighted version must
decide how the mass of a case enters the path. We attach the mass of each case
to its outgoing edge and measure the jumps in weighted ranks.

Formally, remove zero-mass observations, let $m$ be the number remaining, and
relabel them so that $X_1\leq X_2\leq\cdots\leq X_m$. Ties in $X$ are broken
uniformly at random without using $Y$; the response and mass attached to an
observation move with it. Conditional on the realized ordering, define the
weak weighted response rank and survival rank
\begin{align*}
    R_{i,w}
    &= \sum_{j=1}^m p_j\mathbf{1}(Y_j\leq Y_i),
    &
    L_{i,w}
    &= \sum_{j=1}^m p_j\mathbf{1}(Y_j\geq Y_i),
\end{align*}
the weighted path variation and response normalizer
\begin{align*}
    A_w
    &= \sum_{i=1}^{m-1}p_i\left|R_{i+1,w}-R_{i,w}\right|,
    &
    C_w
    &= 2\sum_{i=1}^m p_iL_{i,w}(1-L_{i,w}),
\end{align*}
and the weighted coefficient
\begin{align}
    \widehat\xi_w
    &= 1-\frac{A_w}{C_w}.
    \label{eq:weighted-chatterjee}
\end{align}
The mass $p_i$ multiplies the edge from predictor position $i$ to $i+1$, so a
case contributes to the path variation in proportion to its empirical mass;
the last predictor position has no outgoing edge. An edge weight chosen
independently of the cases would lose the case-weight interpretation of
Section~\ref{sec:framework}.

For uniform masses and an untied response, $R_{i,w}=r_i/m$ with $r_i$ the
ordinary response rank, $A_w = m^{-2}\sum_{i=1}^{m-1}|r_{i+1}-r_i|$, and
$C_w=(m^2-1)/(3m^2)$. Hence \eqref{eq:weighted-chatterjee} reduces exactly to
\begin{align*}
    \widehat\xi
    &= 1-\frac{3\sum_{i=1}^{m-1}|r_{i+1}-r_i|}{m^2-1},
\end{align*}
the continuous-response statistic of \citet{chatterjee2021}. With response
ties and uniform masses, $C_w$ equals Chatterjee's tie-adjusted denominator,
so $C_w$ is its case-weighted extension.

\subsection{Finite-sample properties and computation}

The coefficient is invariant under a common positive rescaling of the weights
and under strictly increasing transformations of either variable. It is
undefined when the response is constant on the positive-mass observations,
because then $C_w=0$; otherwise $\widehat\xi_w\leq1$, and like the ordinary
finite-sample coefficient it can be negative. Because the outgoing edge
carries its base-point mass, reversing the predictor order can change the
coefficient when the masses are unequal; this orientation sensitivity
disappears under uniform masses.

Two conventions require care. Zero-mass observations must be removed before
the predictor is ordered, since retaining one would create a spurious edge
between two positive-mass observations. Predictor ties must be broken without
using the responses: ordering a tied group by its responses would manufacture
short response-rank jumps and hence artificial dependence, whereas uniform tie
breaking preserves the null exchangeability used in
Section~\ref{subsec:chatterjee-test}. A fixed random seed makes the reported
coefficient reproducible; averaging over tie orders would define a different
estimator.

Computation takes $O(m\log m)$ time and $O(m)$ memory: one sort by the
predictor, one sort of the responses for the two weighted ranks, and a linear
scan for $A_w$ and $C_w$. Equation~\eqref{eq:weighted-chatterjee} is the
coefficient reported for both continuous and tied responses. Its independence
test uses a closely related statistic with a deterministic denominator, which
is introduced in Section~\ref{subsec:chatterjee-test}.
